\documentclass[10pt,a4paper]{amsart}
\usepackage[margin=19mm]{geometry}
\usepackage{amsmath,amssymb,mathtools}
\usepackage[scr=rsfs,cal=euler]{mathalfa}
\usepackage{xfrac}
\usepackage[unicode,hidelinks,bookmarks=false]{hyperref}
\newcommand{\ie}{i.e.\ }
\newcommand{\cf}{cf.\ }
\newcommand{\resp}{respectively\ }
\newcommand{\Subsection}{\subsection}
\providecommand{\nobreakdash}{\nobreak\hbox{-}}
\setlength{\emergencystretch}{2em}
\multlinegap0pt
\usepackage{multirow}
\usepackage{amssymb}
\usepackage{caption}
\usepackage{tikz}
\newtheorem{proposition}{Proposition}
\newtheorem{lemma}{Lemma}
\newtheorem{theorem}{Theorem}
\def\C{\mbox{$\mathbb C$}}
\def\CCCC{\mathscr C}
\def\EEE{{\cal E}}
\def\cal{\mathcal}
\def\g{\gamma}
\def\t{\tilde}
\title{Polynomials in molecules}
\author{Yan Gao, Jinsong Zeng}
\date{Source: arXiv:2601.18605v1 (CC BY 4.0)}
\begin{document}
\maketitle

\noindent\emph{Source and licence.} Polynomials in molecules. Version arXiv:2601.18605v1 (CC BY 4.0), distributed by arXiv under the Creative Commons Attribution 4.0 International licence, http://creativecommons.org/licenses/by/4.0/, as recorded in the arXiv licence metadata for this exact version and shown on the abstract page https://arxiv.org/abs/2601.18605v1. Full licence notice: \texttt{licenses/arxiv-2601.18605-CC-BY-4.0.txt}.\par\smallskip

\noindent\textit{Editorial scope.} Counted unit: the article's proposition pro:same (source0.tex lines 1678--1680). The article's complete original proof of it is delivered with it (source0.tex lines 1682--1718, one \texttt{proof} environment of 37 lines). No mathematics is written, repaired or re-derived: the statement and the proof are the article's own lines, in the article's own notation, and no retained line is substituted or re-flowed.  Retained uncounted as the source-local context the proof needs: lines 551--553; lines 705--711; lines 731--749; lines 1655--1662.  One declared adaptation: exactly 1 ancillary strings inside the retained spans are omitted (image-inclusion commands and, where the source repeats a label, the repeated label definitions) and no other line differs from the source; the figure environments, with their placement, captions and labels, are kept, and the package records the omitted strings in fidelity receipts bound to the affected bodies.  No retained line contains a citation, so the wrapper carries no bibliography and no bibliography entry.  The retained lines contain the article's own diagram code, which the wrapper typesets with the standard tikz packages; no image file is delivered and the package declares no asset dependency.  Editorial formatting only, in addition: an AMS \texttt{amsart} preamble that restates the article's own declarations and macros used by the retained lines, namely the environments \texttt{proposition}, \texttt{lemma}, \texttt{theorem} on separate counters, together with the standard packages those lines need.  The retained environments print in this document as Proposition 1, Lemma 1, Theorem 1 in order of appearance, whereas the article prints the same items with the numbering of its own full document.  The complete native source of the article as distributed in the arXiv e-print for this exact version is bundled unchanged as \texttt{sources/193402/source0.tex}.
	\begin{proposition}\label{pro:chain1}
		Suppose that $J_f$ is locally connected. Then the maximal Fatou chain $E_U$ generated by a fixed Fatou domain $U$ has no periodic cut points. In particular, $E_U=\overline{U}$ if and only if $\partial U$ contains no critical points.
	\end{proposition}

\begin{lemma}\label{lem:stable}
Let $\{f_n\}\subseteq\CCCC_d$ be a sequence converging to $g$, and let $z$ be a preperiodic point of $g$.
\begin{enumerate}
\item If $z$ is pre-repelling, then for every $\theta$ with $\pi_g(\theta)=z$, we have  $\overline{R_{f_n}(\theta)}\to\overline{R_g(\theta)}$ as $n\to\infty$ and  $\pi_{f_n}(\theta)$ is pre-repelling for all large $n$. If, in addition, $z$ is not pre-critical, then  $\pi_g(\theta)=\pi_g(\eta)=z$ implies $\pi_{f_n}(\theta)=\pi_{f_n}(\eta)=z_n$ for all sufficiently large $n$.\vspace{2pt}

\item If $z$ is a parabolic periodic point, and if each $f_n$ admits a parabolic point $z_n$ with the same period, multiplier and multiplicity as $z$, such that $z_n\to z$, then for every $\theta$ with $\pi_g(\theta)=z$, we have $\pi_{f_n}(\theta)=z_n$ for all sufficiently large $n$, and $\overline{R_{f_n}(\theta)}\to\overline{R_g(\theta)}$ as $n\to\infty$.
\end{enumerate}\end{lemma}

	\begin{theorem}[Dynamical perturbation]\label{thm:perturbation}
		Let $f\in\CCCC_d$ be  geometrically finite with non-empty bounded Fatou domains. Let $c^*$ be a last critical point of $f$. Then there exists a sequence of geometrically finite maps $\{f_n\}\subseteq \CCCC_d$ uniformly converging to $f$ such that:
		\begin{enumerate}
			\item \textup{(Critical point correspondence)} For every $c\in C_f$,  $f_n$ has a unique critical point $c_n$ with  ${\rm deg}(f,c)={\rm deg}(f_n, c_n)$ and $c_n\to c$ as $n\to\infty$, such that $c_n^*\in F_f$, and  $c_n\not=c_n^*$ lies in a parabolic domain or the boundary of a bounded Fatou domain precisely when $c$ does.\vspace{3pt}
			
			\item \textup{(Critical relation consistency)} For distinct $c,c'\in C_f$, if $f^k(c)=c'$, then $f_n^k(c_n)=c_n'$;
			if $c$ is not iterated to $c^*$, then $f^{m+l}(c)=f^{m}(c)$ implies $f_n^{m+l}(c_{n})=f_n^m(c_{n})$. \vspace{3pt}
			\item\textup{(Parabolic point correspondence)} For each parabolic periodic point $z$ of $f$, there is a unique parabolic point $z_n$ of $f_n$ converging to $z$ as $n\to\infty$, and sharing the same multiplier, period, and multiplicity as $z$ does for $f$.			
			
			\item \textup{(Fatou chain correspondence)} If $c^*$ lies in a maximal Fatou chain of $f$, then there exists a bijection $\iota_n:\EEE_f\to\EEE_{f_n}$ such that\vspace{2pt}
			\begin{itemize}
				\item $c\in K$ if and only if $c_n\in\iota_n(K)$ for any $c\in C_f$ and $K\in\EEE_f$;
				\item $\iota_n\circ f(K)=f_n\circ \iota_n(K)$ for any $K\in\EEE_f$.
			\end{itemize}
			\vspace{3pt}
			
			\item \textup{(Quasiconformal deformation and semiconjugacy)} If $c^*$ lies in the boundary of  a bounded Fatou domain, then there exists a polynomial $\t{f}\in\CCCC_d$ such that each $f_n$ is  quasiconformal conjugate to $\t{f}$. Moreover, there exists a continuous surjection $\phi:J_{\t{f}}\to J_f$ satisfying $$f\circ\phi=\phi\circ \t{f}\ \text{on}\ J_{\t{f}},$$ such that   $\t{z}=\phi^{-1}(z)$ is a singleton and $\pi_f^{-1}(z)=\pi_{\t{f}}^{-1}(\t{z})$ unless $f^k(z) = c^*$ for some $k \ge 0$, in which case $\phi^{-1}(z)$ consists of $d_z=\deg(f^{k+1}, z)$ distinct points $\t{z}_1,\ldots,\t{z}_{d_z}$ satisfying $$\pi_f^{-1}(z)=\bigsqcup_{j=1}^{d_z}\pi_{\t{f}}^{-1}(\t{z}_j).$$
					\end{enumerate}
	\end{theorem}

\begin{lemma}\label{pro:wake}
	Let $f\in\CCCC_d$ be a subhyperbolic map. Let $K$ be the maximal Fatou chain of $f$ generated by a fixed Fatou domain of $f$. Suppose that $W$ is the wake of $K$ bounded by a ray pair $\g$, then for every integer $n\geq 1$,
	\begin{enumerate}
		\item $f^n(\g)$ is  a ray pair that bounds a wake of $K$, denoted by $W_n$;
		\item $f^n$ sends a local sector of $W$ homeomorphically onto a local sector of $W_n$;
		\item the ray pair $\g$ is eventually periodic under iteration by $f$.
	\end{enumerate}
\end{lemma}

\begin{proposition}\label{pro:same}
	Let $f\in\CCCC_d$ be a subhyperbolic map, and $c^*$ be a last critical point of $f$ lying on the boundary of a bounded Fatou domain. Let $g$ be a dynamical perturbation of $(f,c^*)$ as given in Theorem \ref{thm:perturbation}. If every critical point of $f$ is eventually mapped into the maximal Fatou chain generated by a periodic bounded Fatou domain, then the same property still holds for $g$.
\end{proposition}

\begin{proof}
	Assume, by contradiction, that there exists a critical point
	$c_g$ of $g$ that is never iterated into any maximal Fatou chain generated by a bounded Fatou domain.
	
	By Theorem \ref{thm:perturbation}\,(5), there exists a sequence $\{g_n\}$ of polynomials, obtained via quasiconformal deformations of $g$ such that $g_n\to f$ as $n\to\infty$.
	For notational simplicity, if $A_g$ denotes an object associated with $g$ (e.g. an external ray, a preperiodic point, etc.), we write $A_n$ for the corresponding object under $g_n$, and $A:=\lim\limits_{n\to\infty} A_n$ for the corresponding object under $f$.
	
	By the hypothesis of the proposition, there exists a periodic bounded Fatou domain $ U$ of $f$ and an integer $l\geq0$ such that $f^l(c)$ belongs to the maximal Fatou chain $E$ generated by $ U$. Without loss of generality, we may assume that $f( U)= U$.
	
	Let $ U_g$ be the attracting Fatou domain of $g$ corresponding to $ U$, and let $E_g$ be the maximal Fatou chain generated by $ U_g$. Fix an external ray $R_g(\alpha)$ lands at $g^l(c_g)$. By assumption, $g^l(c_g)\notin E_g$. So  $\overline{R_g(\alpha)}$ is contained in a wake $W_g$ of $E_g$ bounded by a ray pair $\g_g=R_g(\theta)\cup \{z_g\}\cup R_g(\theta')$. Since $\g_g$ is $g$-preperiodic by Lemma \ref{pro:wake}, we may assume that
	$\g_g':=g(\g_g)$ is fixed by $g$. Then $w_g:=g(z_g)\in E_g$ is a repelling fixed point of $g$. Denote by $W_g'$ the wake of $E_g$ bounded by $\g'_g$. %See Figure \ref{fig:12}.
%	\begin{figure}
%		\centering
%		
%		\caption{The wake $W_g$ of $E_g$ contains $g^\ell(c_g)$, where $d_0={\rm deg}(g,z_g)=2$.}\label{fig:12}
%	\end{figure}
	
	
	By Lemma \ref{lem:stable}\,(1), $\g_n\to \g=R_f(\theta)\cup\{z\}\cup R_f(\theta')$ as $n\to\infty$. Then $\g$ separates $R_f(\alpha)$ from $ U$. We now verify the following two properties:
	\begin{enumerate}
		\item $f$ sends a local sector of $W$ homeomorphically onto a local sector of $W'$;
		\item $\overline{R_f(\alpha)}$ and $ U$ lie in distinct components of $\mathbb{C}\setminus\g.$
	\end{enumerate}
	
	To show (1), let $d_0:={\rm deg}(g, z_g)\geq 1$. Then there exists $d_0$ external rays $$R_g(\theta)=R_g(\theta_1),\ldots, R_g(\theta_{d_0})$$
	landing at $ z_g$, such that $g$ maps each of them onto $R_g(\sigma_d(\theta))$. Since $g$ maps a local sector of $W_g$ homeomorphically to that of $W_g'$, the wake $W_g$ lies entirely in one component of $$\C\setminus\bigcup_{j=1}^{d_0}\overline{R_g(\theta_j)}.$$
	
	By Theorem \ref{thm:perturbation}, ${\rm deg}(f,  z)=d_0$. Combined with Lemma \ref{lem:stable}, the $d_0$ rays $R_f(\theta_j)$ are all preimages under $f$ of $R_f(\sigma_d(\theta))$ that land at $z$. Since $W$ lies in a component of $\C\setminus\bigcup_{j=1}^{d_0}\overline{R_f(\theta_j)}$, property (1) holds.\vspace{2pt}
	
	To show (2), it suffices to prove that $z$ is not the landing point of $R_f(\alpha)$. Suppose otherwise; then $R_f(\sigma_d(\alpha))$ lands at the repelling fixed point $w$ of $f$. Here $w=\lim_{n\to\infty} w_n$ and $\{w\}=\gamma'\cap J_f$. Since $R_f(\sigma_d(\alpha))\subseteq W'$, it follows from Lemma \ref{lem:stable} that $R_g(\sigma_d(\alpha))\subseteq W_g'$ and  it lands at $w_g$. Hence, there exists $\beta\neq \alpha$ such that $\sigma_d(\beta)=\sigma_d(\alpha)$ and $R_g(\beta)\subseteq W_g$ lands at $z_g$.  Consequently, $R_f(\beta)\subseteq W$ lands at $z$.
	Then the two distinct rays $R_f(\alpha)$ and $R_f(\beta)$ in $W$  are both mapped onto $R_f(\sigma_d(\alpha))$ by $f$, contradicting property (1).\vspace{2pt}
	
	
	We now claim that the fixed point $w$ is a cut point of the maximal Fatou chain $E$ generated by $U$. Indeed,  note first that $R_f(\alpha)$ lands at $f^l(c)$. Since both $f^l(c)$ and $ U$ are contained in $E$, by properties (1) and (2), we may choose a point $z_*\in E\cap W$ in a local sector of $W$ such that $f(z_*)\in E\cap W'$. Then $\g'$ separates $f(z_*)$ from $ U$. So $w$ disconnects $E$.
	
	This claim contradicts Proposition \ref{pro:chain1}. So the  proposition is proved.
\end{proof}

\end{document}
